\section*{Capítulo \ref{ch:g9:heredity-and-traits} --- La herencia y los rasgos}

\begin{solution}{exo:g9:heredity-and-traits:1}
Un \omterm{def:g9:heredity-and-traits:traits}{rasgo} es cualquier característica describible de un individuo; la
\omterm{def:g9:heredity-and-traits:heredity}{herencia} es la transmisión de \omterm{def:g9:heredity-and-traits:traits}{rasgos} de padres a hijos por la
reproducción. La prueba que separa: lo que dos \omterm{def:g8:sexual-reproduction:gametes}{gametos} no pueden
llevar, las familias no lo pueden transmitir.
\end{solution}

\begin{solution}{exo:g9:heredity-and-traits:2}
Grupo sanguíneo: hereditario. \omterm{def:g1:the-five-senses:sense}{Oreja} agujereada: adquirido. Pecas:
tendencia hereditaria, vivida al sol (las dos obras). Francés con
soltura: adquirido. Forma del lóbulo: hereditario.
\end{solution}

\begin{solution}{exo:g9:heredity-and-traits:3}
Patrón familiar (reaparece siguiendo las líneas de la reproducción, o
sigue a las vidas); presencia desde el principio (fijado en la
\omterm{def:g5:flower-to-fruit:steps}{fecundación}, o construido después); transmisión (pasa a los hijos, o
muere con quien lo adquirió).
\end{solution}

\begin{solution}{exo:g9:heredity-and-traits:4}
Porque el \omterm{def:g3:bones-and-muscles:muscle}{músculo} entrenado es adquirido --- construido con el uso, no
llevado por los \omterm{def:g8:sexual-reproduction:gametes}{gametos}. Demuestra que lo que una vida le añade a un
cuerpo no queda escrito en lo que transmite la reproducción.
\end{solution}

\begin{solution}{exo:g9:heredity-and-traits:5}
Transmitir lo que no muestra: el determinante del \omterm{def:g9:heredity-and-traits:traits}{rasgo} viajó
escondido por él y salió a la superficie una generación después.
\end{solution}

\begin{solution}{exo:g9:heredity-and-traits:6}
Porque los \omterm{def:g9:heredity-and-traits:traits}{rasgos} en sí (una barbilla, un grupo sanguíneo) no pueden
viajar en dos \omterm{def:g5:first-look-at-cells:cell}{células} sueltas --- lo que viaja tienen que ser
instrucciones para construirlos.
\end{solution}

\begin{solution}{exo:g9:heredity-and-traits:7}
Heredado: el margen de estatura que sugiere la constitución --- el
nieto recibió los determinantes de la familia. Vivido: una alimentación
más rica rellenó más ese margen --- lo que el ambiente le añade a la
\omterm{def:g9:heredity-and-traits:heredity}{herencia}.
\end{solution}

\begin{solution}{exo:g9:heredity-and-traits:8}
Cada hijo recibe una selección distinta de los determinantes de cada
progenitor --- el barajado reparte combinaciones nuevas en cada
\omterm{def:g5:flower-to-fruit:steps}{fecundación}. Los gemelos idénticos se separan \emph{después} de una
sola \omterm{def:g5:flower-to-fruit:steps}{fecundación}: un reparto, dos construcciones --- así que sus juegos
de determinantes coinciden.
\end{solution}

\begin{solution}{exo:g9:heredity-and-traits:9}
(a) Sigue al árbol, aparece a la edad familiar sin importar el oficio
ni el pueblo, y es concebiblemente llevable por \omterm{def:g8:sexual-reproduction:gametes}{gametos}: hereditario.
(b) Construido con el entrenamiento, sigue al deporte y no a los
padres, y ningún \omterm{def:g8:sexual-reproduction:gametes}{gameto} lleva fondo físico: adquirido --- sobre un
margen heredado de constitución y de \omterm{def:g4:heart-and-blood:heart}{corazón} (la
\cref{prop:g7:organs-at-work:training} construye lo que la \omterm{def:g9:heredity-and-traits:heredity}{herencia}
solo esboza).
\end{solution}

\begin{solution}{exo:g9:heredity-and-traits:10}
Nunca la destreza --- las \omterm{def:g8:nervous-communication:synapse}{sinapsis} se entrenan de nuevo en cada vida.
Lo que sí puede pasar: la materia prima de la aptitud --- oído,
soltura de manos, carácter --- un margen heredado que las casas de
músicos rellenan después con práctica, clases y ejemplo (el segundo
canal de transmisión, el de vivir juntos).
\end{solution}

\begin{solution}{exo:g9:heredity-and-traits:11}
Saltarse una generación es la \omterm{def:g9:heredity-and-traits:heredity}{herencia} comportándose con normalidad:
cada progenitor puede llevar sin mostrarlo el determinante de los \omterm{def:g1:the-five-senses:sense}{ojos}
azules y pasarlo; dos de esas copias escondidas que se encuentran en un
hijo enseñan el \omterm{def:g9:heredity-and-traits:traits}{rasgo}. El patrón demuestra que se lleva sin mostrar ---
y nada sobre la honradez de la familia.
\end{solution}

\begin{solution}{exo:g9:heredity-and-traits:12}
Todo lo que transmite el padre viaja en el \omterm{def:g8:reproductive-systems:male}{espermatozoide}, y el
\omterm{def:g8:reproductive-systems:male}{espermatozoide} es una cola y un \omterm{def:g6:cell-unit-of-life:parts}{núcleo} apretado --- así que los
determinantes tienen que estar en el \omterm{def:g6:cell-unit-of-life:parts}{núcleo}. Y como las instrucciones
paternas de un ser humano entero caben en un paquete de milésimas de
milímetro de ancho, las instrucciones tienen que estar escritas a
pequeñez molecular --- una letra más fina que ninguna estructura
celular que enseñe el \omterm{def:g5:first-look-at-cells:microscope}{microscopio} del colegio.
\end{solution}

\begin{solution}{pb:g9:heredity-and-traits:1}
\textbf{1.} Barbilla de la capilla: hereditaria (patrón familiar,
presente desde el principio). Ser zurdo: tendencia hereditaria (sigue a
las familias sin rigor, fijada temprano). Tos de los panaderos:
adquirida (sigue al oficio). Grupos sanguíneos: \omterm{def:g9:heredity-and-traits:heredity}{hereditarios} (fijados
en la \omterm{def:g5:flower-to-fruit:steps}{fecundación}, llevables por \omterm{def:g8:sexual-reproduction:gametes}{gametos}). Caras curtidas: adquiridas
(siguen a la vida de pastor).

\textbf{2.} La prueba del patrón familiar: el \omterm{def:g9:heredity-and-traits:traits}{rasgo} se dibuja sobre las
líneas de un oficio, no sobre las de la reproducción --- casarse con
alguien de la casa lo adquiere, mudarse lejos lo quita, y no interviene
ningún \omterm{def:g8:sexual-reproduction:gametes}{gameto}.

\textbf{3.} Que un determinante se puede llevar sin enseñarlo: la
generación de en medio transmitió las instrucciones de la barbilla
mientras enseñaba otras barbillas.

\textbf{4.} Cada progenitor llevaba, junto al determinante del grupo
que enseñaba, un determinante del grupo O escondido --- y el hijo
recibió los dos escondidos. Se lleva sin mostrar, con la tinta del
propio registro.

\textbf{5.} No la barbilla, sino sus determinantes --- instrucciones
para construirla --- viajando en los únicos vehículos que tiene la
reproducción: un \omterm{def:g4:animal-reproduction:sexual}{óvulo} y un \omterm{def:g8:reproductive-systems:male}{espermatozoide} por generación, cuatro
entregas en el siglo.

\textbf{6.} En las selecciones: cada \omterm{def:g5:flower-to-fruit:steps}{fecundación} repartió el
determinante compartido de la barbilla junto a una selección nueva a
medias de todo lo demás --- carta común, manos distintas.

\textbf{7.} Los \omterm{def:g9:heredity-and-traits:traits}{rasgos} adquiridos no se transmiten nunca. Razón: se
construyen en los tejidos del cuerpo al vivir, y los \omterm{def:g8:sexual-reproduction:gametes}{gametos} --- dos
\omterm{def:g5:first-look-at-cells:cell}{células}, que llevan solo determinantes --- no tienen manera de
anotarlos.

\textbf{8.} Una \omterm{def:g5:flower-to-fruit:steps}{fecundación} y una separación posterior: un solo reparto
de determinantes, construido dos veces --- los gemelos idénticos de la
\cref{rem:g5:human-reproduction:twins}, certificados por las
coincidencias del libro.

\textbf{9.} Canal de los \omterm{def:g8:sexual-reproduction:gametes}{gametos}: la barbilla, los grupos sanguíneos,
la tendencia a ser zurdo. Canal de vivir juntos: la tos (el aire de la
panadería), el curtido de la cara (el tiempo de los montes) --- y, del
capítulo del álbum, las lenguas, los oficios y las destrezas.

\textbf{10.} Por los saltos: la generación de en medio llevaba la
barbilla sin enseñarla. Por el registro de \omterm{prop:g4:journey-of-food:absorb}{sangre}: unos padres que
enseñaban A y B llevaban O sin enseñarlo. Dos veces, lo que se lleva
supera a lo que se enseña.

\textbf{11.} Físicamente: están en los \omterm{def:g6:cell-unit-of-life:parts}{núcleos} de los \omterm{def:g8:sexual-reproduction:gametes}{gametos} --- eso
ya lo demuestra la anatomía del \omterm{def:g8:reproductive-systems:male}{espermatozoide}. De qué están hechas, y
cómo escribe un paquete tan pequeño un libro tan grande, es
justamente la pregunta de los capítulos siguientes.

\textbf{12.} La mayoría de los \omterm{def:g9:heredity-and-traits:traits}{rasgos} son márgenes heredados, vividos
hasta el fondo: los \omterm{def:g8:sexual-reproduction:gametes}{gametos} reparten el boceto, la vida lo dibuja.
\end{solution}
